\documentclass[10pt,a4paper]{amsart}
\usepackage[margin=19mm]{geometry}
\usepackage{amsmath,amssymb,mathtools}
\usepackage[scr=rsfs,cal=euler]{mathalfa}
\usepackage{xfrac}
\usepackage[unicode,hidelinks,bookmarks=false]{hyperref}
\newcommand{\ie}{i.e.\ }
\newcommand{\cf}{cf.\ }
\newcommand{\resp}{respectively\ }
\newcommand{\Subsection}{\subsection}
\providecommand{\nobreakdash}{\nobreak\hbox{-}}
\setlength{\emergencystretch}{2em}
\multlinegap0pt
\usepackage{booktabs}
\usepackage{amssymb}
\usepackage{mathtools}
\usepackage{float}
\usepackage{natbib}
\usepackage{xcolor}
\usepackage[shortlabels]{enumitem}
\newtheorem{proposition}{Proposition}
\usepackage{cleveref}
\crefname{proposition}{Proposition}{Propositions}
\newcommand{\PP}{\mathbb{P}}
\newcommand{\R}{\mathbb{R}}
\newcommand{\Z}{\mathbb{Z}}
\newcommand{\eps}{\varepsilon}
\newcommand{\sign}{\operatorname{sign}}
\newcommand{\stepsize}{\alpha}
\title{Insights on Muon from Simple Quadratics}
\author{Antoine Gonon \and Andreea-Alexandra Muşat \and Nicolas Boumal \and \small Institute of Mathematics, EPFL, Switzerland}
\date{Source: arXiv:2602.11948v1 (CC BY 4.0)}
\begin{document}
\maketitle

\noindent\emph{Source and licence.} Insights on Muon from Simple Quadratics. Version arXiv:2602.11948v1 (CC BY 4.0), distributed by arXiv under the Creative Commons Attribution 4.0 International licence, http://creativecommons.org/licenses/by/4.0/, as recorded in the arXiv licence metadata for this exact version and shown on the abstract page https://arxiv.org/abs/2602.11948v1. Full licence notice: \texttt{licenses/arxiv-2602.11948-CC-BY-4.0.txt}.\par\smallskip

\noindent\textit{Editorial scope.} Counted unit: the article's proposition prop:noise-breaks-grid (source0.tex lines 1799--1825). The article's complete original proof of it is delivered with it (source0.tex lines 1827--1837, one \texttt{proof} environment of 11 lines). No mathematics is written, repaired or re-derived: the statement and the proof are the article's own lines, in the article's own notation, and no retained line is substituted or re-flowed.  Retained uncounted as the source-local context the proof needs: lines 1881--1881.  Nothing was omitted from any retained span, so the package carries no omission record.  No retained line contains a citation, so the wrapper carries no bibliography and no bibliography entry.  No retained line contains a figure environment, an image-inclusion command or drawing code, so no image is delivered and the package declares no asset dependency.  Editorial formatting only, in addition: an AMS \texttt{amsart} preamble that restates the article's own declarations and macros used by the retained lines, namely the environments \texttt{proposition} on separate counters, together with the standard packages those lines need.  The retained environments print in this document as Proposition 1 in order of appearance, whereas the article prints the same items with the numbering of its own full document.  The complete native source of the article as distributed in the arXiv e-print for this exact version is bundled unchanged as \texttt{sources/194506/source0.tex}.
\subsection{Proof of \Cref{prop:noise-breaks-grid}} \label{sec:noise-breaks-grid}

\begin{proposition}[Noise breaks the deterministic grid trap]
\label{prop:noise-breaks-grid}
Consider the one-dimensional noisy sign dynamics with step size $\stepsize>0$
and noise level $\sigma\ge 0$:
\[
  w_{n+1}^{(\sigma)}
  = w_n^{(\sigma)} - \stepsize\bigl(\sign(w_n^{(\sigma)}) + \sigma \xi_{n+1}\bigr),
  \qquad n \ge 0,
\]
where $(\xi_n)_{n\ge1}$ are i.i.d.\ $\mathcal{N}(0,1)$ random variables and $w_0^{(\sigma)}\in\R$ is fixed. 

Fix $\eps>0$ and define the hitting time
\begin{align}
  T_\eps^{(\sigma)} := \inf\{n \ge 0 : |w_n^{(\sigma)}| \le \eps\}.
  \label{eq:hittingtimefoo}
\end{align}

Assume $\sigma>0$. 
Then, for any initialization $w_0\in\R$,
\[
  \PP\bigl(T_\eps^{(\sigma)} < \infty\bigr) = 1.
\]

By contrast, for the deterministic dynamics ($\sigma = 0$), if the grid
$w_0 + \stepsize\Z$ does not intersect $[-\eps,\eps]$, then
$T_\eps^{(0)} = \infty$ for all such initializations.
\end{proposition}

\begin{proof}[Proof idea for \Cref{prop:noise-breaks-grid}]
  (Full proof in \Cref{sec:noise-breaks-grid}.)
The proof combines a drift-to-compact argument
with a uniform ``attempt'' mechanism:
once the process enters a fixed neighborhood of the origin, 
there is a state-uniform positive probability to land in $[-\eps,\eps]$ 
within a bounded number of steps, 
and repeated returns to that neighborhood 
(which happen almost surely because of the drift)
guarantee eventual success.
\end{proof}

\end{document}
